\begin{figure}[t]
  \centering
  % Laid out in explicit millimetres and scaled to the column. Node *text* does not scale with the
  % x/y unit vectors, so every width below is measured with \settowidth rather than assumed: the
  % example line is 59.6mm and the detector box's widest line 32.8mm at \scriptsize. Assuming them
  % is what put the sentences underneath the measurement boxes, and the box text outside its box.
  \resizebox{\textwidth}{!}{%
  \begin{tikzpicture}[
      x=1mm,y=1mm,
      box/.style={draw=blue!55!black, fill=blue!6, rounded corners=1.2mm,
                  line width=.45pt, align=center, inner sep=1.5mm, font=\scriptsize},
      measure/.style={box, draw=orange!70!black, fill=orange!8},
      line/.style={align=left, font=\scriptsize, inner sep=.5mm},
      tag/.style={font=\scriptsize\itshape, text=black!55, inner sep=.5mm},
      hit/.style={fill=teal!16, rounded corners=.6mm, inner sep=.5mm},
      kept/.style={fill=red!12, rounded corners=.6mm, inner sep=.5mm},
      flow/.style={-{Latex[length=1.5mm]}, line width=.45pt, draw=black!70},
      rail/.style={line width=.45pt, draw=black!70},
      rewrite/.style={-{Latex[length=1.6mm]}, line width=.5pt, draw=teal!55!black, dashed}
    ]
    \node[box, minimum width=37mm, minimum height=16mm] (det) at (18.5,30)
      {detector pool\\\scriptsize 15 rules, encoders, LLMs\\[-.3mm]
       \scriptsize union / vote / intersection};

    % One rail out of the detector pool, then a short arrow into each condition: three separate
    % arrows from the box corner converge on the tags and land on top of them.
    \draw[flow] (det.east) -- (41,30);
    \draw[rail] (41,22) -- (41,38);
    \foreach \y in {38,30,22} { \draw[flow] (41,\y) -- (44.5,\y); }

    \node[tag, anchor=east] at (50,38) {A};
    \node[line, anchor=west] at (51.5,38)
      {Frau \tikz[baseline]{\node[hit,anchor=base]{Anna Weber};}, geb.\
       \tikz[baseline]{\node[kept,anchor=base]{12.03.1948};},
       \tikz[baseline]{\node[hit,anchor=base]{Klinik Nord};}.};

    \node[tag, anchor=east] at (50,30) {B};
    \node[line, anchor=west] at (51.5,30)
      {Frau \tikz[baseline]{\node[hit,anchor=base]{Marta Keller};}, geb.\
       \tikz[baseline]{\node[kept,anchor=base]{12.03.1948};},
       \tikz[baseline]{\node[hit,anchor=base]{Klinik S\"ud};}.};

    \node[tag, anchor=east] at (50,22) {C};
    \node[line, anchor=west] at (51.5,22)
      {Frau \tikz[baseline]{\node[hit,anchor=base]{\texttt{[PERSON]}};}, geb.\
       \tikz[baseline]{\node[kept,anchor=base]{12.03.1948};},
       \tikz[baseline]{\node[hit,anchor=base]{\texttt{[ORG]}};}.};

    \draw[rewrite] (46.8,28.4) -- (46.8,23.6);
    \node[tag, anchor=west, text=teal!45!black] at (51.5,15)
      {identical spans, so B rewrites to C};

    \node[measure, minimum width=27mm] (m1) at (131,38) {detection\\and exposure};
    \node[measure, minimum width=27mm] (m2) at (131,30) {downstream\\utility};
    \node[measure, minimum width=27mm] (m3) at (131,22) {attacks and\\stability};
    \foreach \y in {38,30,22} { \draw[flow] (116,\y) -- (117.5,\y); }
  \end{tikzpicture}%
  }
  \caption{One set of detector spans generates three release conditions, on an invented header line
  in the style of a discharge letter. Detected identifiers (teal) become surrogates in B and typed
  placeholders in C; dates (red) pass through by design. B and C replace the same spans, so C is a
  rewrite of B. Every condition meets every measurement family on the same documents.}
  \label{fig:overview}
\end{figure}
